The results of this thesis suggest several promising directions for future research.

\textbf{Algorithmic Extension}

One natural next step is to explore alternative graph similarity measures beyond the edge overlap measure used here. Measures such as graph edit distance or maximum common subgraph could provide a more nuanced assessment of graph similarity and may yield different amplification behaviours under NVQWOA dynamics. Investigating how the algorithm performs with cost functions that include structural penalties or weightings could also reveal how sensitive it is to the choice of objective and inform the design of mixers tailored to specific graph properties. 

\textbf{Methodological Adevances}

The present study was limited to small problem sizes by the computational cost of state-vector simulation. Extending the analysis to larger graphs using tensor-network simulation or distributed simulation frameworks would provide  insight into the algorithm's scaling behaviour. Additionally, implementing NVQWOA on noisy intermediate-scale quantum (NISQ) hardware would allow experimental validation of its performance under realistic noise and decoherence conditions. Such tests could help assess the robustness of amplification under hyperparameter selections and inform strategies for error mitigation or noise-aware parameter selection.

\textbf{Incorporating Problem Structure into the Initial State}

In the present formulation, NVQWOA begins from an initial superposition over all possible solutions, which treats all states as equally likely. Although this is easy to implement, it gives up the opportunity to exploit graph structures that can be easily read from any vertex mapping, such as degree distribution or spectral properties. Using these to inform the structure of the initial state could strongly constrain the solution space of plausible vertex mappings, which would reduce the search space that the algorithm explores.